	\subsection{Kiểm định Breusch--Pagan (1979)}
	\textbf{Khi nào cần dùng:} Khi ta có cơ sở nghi ngờ rằng phương sai của sai số thay đổi (lớn dần hoặc nhỏ dần) theo \textit{hàm tuyến tính} của một hoặc nhiều biến độc lập.
	
	\textbf{Cơ sở lý thuyết \& Chứng minh:} Breusch và Pagan \cite{breusch1979} xây dựng kiểm định này dựa trên nguyên lý \textit{Nhân tử Lagrange (Lagrange Multiplier - LM)} trong phương pháp Ước lượng Hợp lý Cực đại (MLE). Giả sử phương sai tuân theo hàm $\sigma_i^2 = h(\gamma_0 + \mathbf{z}_i^{\top}\boldsymbol{\gamma})$. Đạo hàm của hàm log-likelihood (hàm Score) được đánh giá tại $\boldsymbol{\gamma} = \mathbf{0}$. Nếu đạo hàm này khác không một cách có ý nghĩa thống kê, ta có cơ sở bác bỏ tính đồng nhất.
	
	\textbf{Thiết lập giả thuyết:}
	\begin{itemize}
		\item $H_0: \gamma_1 = \gamma_2 = \dots = \gamma_k = 0$ (Phương sai là hằng số / Đồng nhất phương sai).
		\item $H_1:$ Có ít nhất một $\gamma_j \neq 0$ (Phương sai thay đổi).
	\end{itemize}
	
	\textbf{Quy trình và Công thức:}
	\begin{enumerate}
		\item Chạy mô hình OLS gốc, lưu bình phương phần dư: $\hat{\varepsilon}_i^2$.
		\item Chạy hồi quy phụ: $\hat{\varepsilon}_i^2 = \gamma_0 + \gamma_1 X_{i1} + \dots + \gamma_k X_{ik} + v_i$.
		\item Tính thống kê kiểm định: $LM = \frac{1}{2}\text{ESS} \sim \chi^2(k)$. 
	\end{enumerate}
	\textbf{Làm sao để thỏa \& Tối ưu hóa:} Để thỏa mãn giả định OLS, ta cần chấp nhận $H_0$ (tức là $p\text{-value} > \alpha$, thường $\alpha = 0.05$). Tuy nhiên, phiên bản gốc của B-P rất nhạy cảm với giả định sai số phân phối chuẩn. Để tối ưu, Koenker (1981) \cite{koenker1981} đã đề xuất \textbf{Kiểm định B-P dạng Studentized}, tính toán qua $LM = n \times R^2_{\text{aux}}$. Dạng tối ưu này vững (robust) hơn khi sai số có đuôi phân phối dày và được các phần mềm như R (hàm \texttt{bptest}) sử dụng làm mặc định.
	
	\subsection{Kiểm định White (1980)}
	\textbf{Khi nào cần dùng:} Khi người nghiên cứu \textit{không biết trước dạng cấu trúc} của phương sai thay đổi (có thể là phi tuyến, bậc 2, hoặc do sự tương tác chéo giữa các biến) và tập mẫu dữ liệu đủ lớn.
	
	\textbf{Cơ sở lý thuyết \& Chứng minh:} Halbert White \cite{white1980} phát triển kiểm định này dựa trên thuộc tính của \textit{Đẳng thức Ma trận Thông tin (Information Matrix Equality)}. Theo lý thuyết tiệm cận, nếu $H_0$ đúng, trung bình mẫu của ma trận $\hat{\varepsilon}_i^2 \mathbf{x}_i \mathbf{x}_i^\top$ sẽ hội tụ chính xác về $\hat{\sigma}^2 \frac{1}{n}\sum \mathbf{x}_i \mathbf{x}_i^\top$. Để kiểm tra sự hội tụ này một cách thực tiễn mà không cần thao tác ma trận phức tạp, White chứng minh rằng chỉ cần hồi quy $\hat{\varepsilon}_i^2$ lên mọi phần tử duy nhất cấu thành nên $\mathbf{x}_i \mathbf{x}_i^\top$.
	
	\textbf{Thiết lập giả thuyết:} Tương tự B-P, $H_0:$ Phương sai đồng nhất ($\sigma_i^2 = \sigma^2$ với mọi $i$).
	
	\textbf{Quy trình và Công thức:}
	Hồi quy phụ của kiểm định White cực kỳ bao quát:
	\begin{equation}
		\hat{\varepsilon}_i^2 = \gamma_0 + \sum \gamma_j X_j + \sum \delta_j X_j^2 + \sum_{j<l} \theta_{jl} X_j X_l + v_i
	\end{equation}
	Thống kê kiểm định: $W = n \times R^2_{\text{aux}} \sim \chi^2(p)$, với $p$ là tổng số biến giải thích trong mô hình phụ.
	
	\textbf{Làm sao để thỏa \& Tối ưu hóa:} Cần $p\text{-value} > \alpha$ để kết luận mô hình thỏa mãn đồng nhất phương sai. \textit{Tối ưu:} Nhược điểm chí mạng của White test là làm cạn kiệt bậc tự do cực nhanh (ví dụ mô hình gốc 5 biến, mô hình phụ White sẽ có tới 20 biến). Để tối ưu trong mẫu nhỏ, người ta thường dùng \textbf{Kiểm định White rút gọn (Modified White Test)}: chỉ hồi quy $\hat{\varepsilon}_i^2$ theo giá trị dự báo $\hat{y}_i$ và bình phương của nó $\hat{y}_i^2$, giúp tiết kiệm bậc tự do tối đa mà vẫn giữ được sức mạnh kiểm định.
